% =====================================================================
%  CHAPTER 2: 影响生长发育的因素（对应大纲第二章）
%  内容来源：往届笔记第六章，参考PPT第六章
% =====================================================================
\chapter{影响生长发育的因素}
\setcounter{chapter}{2}
\chapterintro{
\textbf{掌握：}影响生长发育的遗传因素、环境因素，促进生长发育的措施。生长发育长期变化的概念。\newline
\textbf{熟悉：}生长发育长期变化的主要表现。\newline
\textbf{其他：}生长发育长期变化的变化原因及其对人类的影响。
}

% ===== 第一部分：掌握内容 =====
\section{影响生长发育的遗传因素}\difficulty

\subsection{家族性遗传影响}

\begin{defbox}
\textbf{遗传（heredity）：}形态结构、生理功能及身心发育等性状在亲代与子代间的相似性和连续性。
\end{defbox}

\begin{keybox}
\textbf{1. 体格发育的家族性遗传}
\begin{enumerate}[leftmargin=*]
    \item 生长突增模式、性成熟早晚及月经初潮年龄等与家族遗传密切相关。
    \item 成年身高与父母平均身高间的遗传度为0.75，即身高75\%取决于遗传因素。
    \item \textbf{生长发育的家族聚集性（familial aggregation）：}父母与子女身高的相关系数呈随年龄上升的趋势，提示遗传因素在个体越接近成熟的阶段表现得越充分的现象。
\end{enumerate}

\textbf{2. 心理行为发展的遗传影响}
\begin{enumerate}[leftmargin=*]
    \item 遗传对感知觉、气质有较直接的影响。
    \item 在个性品质、道德行为方面，遗传因素对心理-行为的影响作用随年龄增大而减弱，尤其在青少年阶段，遗传因素的作用远不如环境、教育等因素的影响显著。
    \item 环境对某种心理特性或行为的发展所起的作用，往往有赖于这种特性或行为的遗传基础。
\end{enumerate}
\end{keybox}

\subsection{表观遗传}

\begin{defbox}
\textbf{表观遗传（epigenetics）：}主要研究不涉及基因核苷酸序列改变的基因表达和调控的可遗传修饰。其遗传方式有DNA序列不变而表型改变，改变有可遗传性和可逆性，不遵循孟德尔遗传定律等特点。表观遗传是环境因素和细胞内的遗传物质间发生交互作用的结果，使生物在环境不断变化的情况下，能维持遗传序列的稳定。
\end{defbox}

\section{影响生长发育的环境因素}\difficulty

\subsection{物质环境因素}

\begin{keybox}
\textbf{一、自然地理环境因素}
\begin{enumerate}[leftmargin=*]
    \item 气候因素
    \item 季节因素：春季身高增长快，秋季体重增长快。
\end{enumerate}

\textbf{二、环境污染因素}

\textbf{1. 化学性环境污染因素}
\begin{enumerate}[leftmargin=*]
    \item 空气污染
    \item 铅污染
    \item \textbf{环境内分泌干扰物（environmental endocrine disrupters）：}对机体内天然激素的产生、释放、运输、代谢、消除、结合、功能发挥产生干扰作用，从而破坏机体内环境平衡稳定状态的维持，并对个体机能和发育造成影响的外源性环境化学物质。
    \begin{enumerate}[leftmargin=*]
        \item 分类：模拟/干扰雌激素、干扰睾酮、干扰甲状腺素、干扰其他内分泌功能。
        \item 人体暴露途径：消化道、呼吸道和皮肤接触。
        \item 对儿童青少年身心健康的影响：
        \begin{itemize}
            \item 生殖毒性，生命早期敏感窗口期暴露，影响子代生殖器官发育和精子成熟。
            \item 性发育异常，是儿童性早熟的直接病因或促进因素。
            \item 神经毒性，引起持久、不可逆的学习能力缺失和行为发育障碍。
            \item 抑制免疫反应。
        \end{itemize}
    \end{enumerate}
\end{enumerate}

\textbf{2. 物理性环境污染因素：}噪声污染、电磁辐射污染、放射性污染。
\end{keybox}

\subsection{社会决定因素}

\begin{keybox}
\textbf{一、社会经济和医疗卫生保健：}社会经济状况、医疗卫生保障。

\textbf{二、家庭生活质量：}家庭经济状况、父母受教育程度、家庭结构、家庭氛围、亲子关系、教养方式。

\textbf{三、文化和教育}

\textbf{四、现代媒体和娱乐方式：}电视、网络使用、玩具和读物。
\end{keybox}

\subsection{行为生活方式因素}

\begin{defbox}
\textbf{生活方式（life styles）：}人们在一定社会文化、经济、风俗、家庭等因素影响下形成的一系列日常生活习惯和生活模式。

\textbf{行为（behaviors）：}表现为其具体外显，主要包括饮食（早中晚3:4:3）、运动（体育锻炼）、睡眠、娱乐、消费、社会交往等。
\end{defbox}

\section{促进生长发育的措施}

\begin{tipbox}
综合影响生长发育的各类因素，促进儿童少年生长发育的主要措施包括：

\begin{enumerate}[leftmargin=*]
    \item \textbf{优化遗传咨询与孕期保健}——从生命早期开始关注遗传和环境因素的交互作用。
    \item \textbf{改善物质环境}——减少环境污染暴露，特别是环境内分泌干扰物的暴露；利用季节因素（春季促进身高增长，秋季关注体重增长）。
    \item \textbf{优化社会决定因素}——改善家庭生活质量，提高父母受教育程度，营造良好家庭氛围和亲子关系。
    \item \textbf{培养健康行为生活方式}——合理饮食（三餐热量分配3:4:3）、充足运动、规律睡眠。
    \item \textbf{加强学校卫生服务}——定期监测生长发育，早期发现和干预生长发育偏离。
    \item \textbf{关注敏感期}——在生长发育关键窗口期（如生命早期1000天）提供最优环境和支持。
\end{enumerate}
\end{tipbox}

% ===== 第二部分：熟悉内容 =====
\section{生长发育长期变化的主要表现}

\begin{keybox}
\textbf{生长发育长期变化（secular trend of growth）：}指儿童少年生长发育水平随年代演进而持续增长的趋势，是过去一个多世纪以来人类最重要的生物学现象之一。

\textbf{主要表现：}

\begin{enumerate}[leftmargin=*]
    \item \textbf{身高、体重增长}——各年龄组儿童少年的身高、体重普遍出现年代性增长，成年身高逐步提高。
    \item \textbf{青春发育提前}——月经初潮年龄和首次遗精年龄提前，青春发动时相整体前移。
    \item \textbf{体型变化}——身体比例发生变化，体格发育呈现向更高更重的方向发展。
    \item \textbf{牙齿发育提前}——恒牙萌出时间提前。
    \item \textbf{骨龄提前}——骨骼成熟速度加快。
\end{enumerate}
\end{keybox}

% ===== 第三部分：其他内容 =====
\section{生长发育长期变化的原因及其对人类的影响}

\begin{tipbox}
\textbf{变化原因：}

\begin{enumerate}[leftmargin=*]
    \item \textbf{营养改善}——蛋白质和热量摄入增加，营养状况普遍改善是最主要因素。
    \item \textbf{疾病控制}——传染病发病率下降，儿童患病率降低。
    \item \textbf{社会经济条件改善}——生活水平提高，医疗卫生条件改善。
    \item \textbf{异质远缘婚配增加}——遗传多样性增加产生杂种优势效应。
    \item \textbf{城市化进程}——更多刺激和更好的生活条件。
\end{enumerate}

\textbf{对人类的影响：}

\begin{enumerate}[leftmargin=*]
    \item \textbf{积极影响}：体格发育水平提高，人口素质增强；青春期提前意味着更早获得社会功能。
    \item \textbf{消极影响}：体格发育与社会心理发育不协调；青春发动时相提前带来更多心理卫生问题和健康风险（如过早性行为、青少年妊娠）；身高体重增长增加资源消耗和环境压力。
    \item \textbf{卫生学应对}：需要定期更新生长发育参考标准；加强青春期健康教育；关注青春发动时相异常。
\end{enumerate}
\end{tipbox}

% ----- 复习题 -----
\reviewcontent{
\section{复习题}

\begin{enumerate}[leftmargin=*, itemsep=8pt]
    \item \textbf{【名词解释】}遗传；表观遗传；生长发育的家族聚集性；环境内分泌干扰物；生活方式；行为；生长发育长期变化
    \item \textbf{【简答题】}简述影响生长发育的遗传因素。
    \item \textbf{【简答题】}简述环境内分泌干扰物（EDCs）对儿童少年健康的影响。
    \item \textbf{【简答题】}简述影响生长发育的社会决定因素。
    \item \textbf{【简答题】}简述促进儿童少年生长发育的主要措施。
    \item \textbf{【简答题】}简述生长发育长期变化的主要表现及其卫生学意义。
\end{enumerate}

\subsection*{参考答案要点}

\begin{enumerate}[leftmargin=*, itemsep=5pt]
    \item \textbf{遗传}：形态结构、生理功能及身心发育等性状在亲代与子代间的相似性和连续性。\textbf{表观遗传}：研究不涉及基因核苷酸序列改变的基因表达和调控的可遗传修饰，DNA序列不变而表型改变，有可遗传性和可逆性，不遵循孟德尔遗传定律。\textbf{生长发育的家族聚集性}：父母与子女身高的相关系数呈随年龄上升的趋势，提示遗传因素在个体越接近成熟的阶段表现得越充分的现象。\textbf{环境内分泌干扰物}：对机体内天然激素的产生、释放、运输、代谢、消除、结合、功能发挥产生干扰作用，从而破坏机体内环境平衡稳定状态的维持，并对个体机能和发育造成影响的外源性环境化学物质。\textbf{生活方式}：人们在一定社会文化、经济、风俗、家庭等因素影响下形成的一系列日常生活习惯和生活模式。\textbf{行为}：表现为其具体外显，主要包括饮食（早中晚3:4:3）、运动（体育锻炼）、睡眠、娱乐、消费、社会交往等。\textbf{生长发育长期变化}：儿童少年生长发育水平随年代演进而持续增长的趋势，是过去一个多世纪以来人类最重要的生物学现象之一。

    \item (1) 家族性遗传影响：身高75\%取决于遗传因素，生长发育的家族聚集性提示遗传因素在越接近成熟的阶段表现越充分；(2) 表观遗传：环境因素和遗传物质间发生交互作用的结果，DNA序列不变而表型改变，有可遗传性和可逆性；(3) 遗传对心理行为的影响：对感知觉、气质有较直接影响，对个性品质、道德行为的影响随年龄增大而减弱。

    \item (1) 生殖毒性：生命早期敏感窗口期暴露影响子代生殖器官发育和精子成熟；(2) 性发育异常：是儿童性早熟的直接病因或促进因素；(3) 神经毒性：引起持久、不可逆的学习能力缺失和行为发育障碍；(4) 抑制免疫反应。

    \item (1) 社会经济和医疗卫生保健；(2) 家庭生活质量：家庭经济状况、父母受教育程度、家庭结构、家庭氛围、亲子关系、教养方式；(3) 文化和教育；(4) 现代媒体和娱乐方式。

    \item (1) 优化遗传咨询与孕期保健；(2) 改善物质环境，减少环境污染暴露；(3) 优化社会决定因素，改善家庭生活质量；(4) 培养健康行为生活方式（合理饮食、充足运动、规律睡眠）；(5) 加强学校卫生服务，定期监测生长发育；(6) 关注敏感期，在关键窗口期提供最优环境。

    \item 主要表现：(1)身高体重增长，(2)青春发育提前，(3)体型变化，(4)牙齿发育提前，(5)骨龄提前。卫生学意义：需定期更新生长发育参考标准，加强青春期健康教育，关注青春发动时相异常。
\end{enumerate}
}

\newpage
